% Part: first-order-logic
% Chapter: introduction
% Section: soundness-completeness

\documentclass[../../../include/open-logic-section]{subfiles}

\begin{document}

\olfileid{fol}{int}{scp}

\olsection{可靠性与完备性}

我们还将为一阶逻辑引入推导系统。逻辑学家已经发展出许多种推导系统，但它们都定义了语句之间相同的可推导关系。如果存在某种精确定义类型的推导，我们就说 $\Gamma$ \emph{推导}~$!A$，记作 $\Gamma \Proves !A$。推导总是符号的有限排列——可能是一列语句，也可能是某种更复杂的结构。推导系统的目的是提供一种工具，用以判定某个语句是否为某个集合~$\Gamma$ 的语义后承。为了达到这个目的，就必须满足 $\Gamma \Entails !A$ \emph{当且仅当} $\Gamma \Proves !A$。

如果 $\Gamma \Proves !A$ 但并非 $\Gamma \Entails !A$，那么我们的推导系统就太强了，推出了太多东西。如果 $\Gamma \Proves !A$ 则 $\Gamma \Entails !A$，这一性质称为\emph{可靠性}，它是任何一个好的推导系统最起码的要求。另一方面，如果 $\Gamma \Entails !A$ 但并非 $\Gamma \Proves !A$，那么我们的推导系统就太弱了，证明得不够。如果 $\Gamma \Entails !A$ 则 $\Gamma \Proves !A$，这一性质称为\emph{完备性}。可靠性通常相对容易证明（通过对归纳定义的推导结构施归纳）。完备性则更难证明。

可靠性与完备性引出了一系列重要推论。如果一个语句集合~$\Gamma$ 能推导出一个矛盾（例如 $!A \land \lnot !A$），就称它是\emph{不一致的}。不一致的 $\Gamma$ 不可能有任何模型，它们是不可满足的。从完备性可以得到其逆命题：任何并非不一致的——或者如我们之后所称的，\emph{一致的}——$\Gamma$ 都有一个模型。事实上，这与完备性是等价的，并且是我们将要实际证明的完备性形式。这是一个深刻而或许令人惊讶的结果：仅仅由于无法从 $\Gamma$ 证明 $!A \land \lnot !A$，就能保证存在一个如 $\Gamma$ 所描述的结构。所以，完备性回答了这样一个问题：哪些语句集合有模型？答案是：所有且仅有一致的集合。

可靠性与完备性定理还有两个重要推论：紧致性定理和 Löwenheim--Skolem定理。这些是模型论中的重要结果，可以用来确立许多有趣的结论。我们已经提到过两个：一阶逻辑无法表达一个结构的论域是有穷的，也无法表达它是不可数的。

从历史上看，所有这些工作——如何定义一阶逻辑的语法和语义、如何定义好的推导系统、如何证明它们可靠且完备、厘清一阶语言能和不能表达什么——都花了很长时间才得以澄清和完善。我们现在已然知晓其法，但逐一梳理所有细节仍可能令人困惑与繁琐。但这也很重要，因为这里为一阶逻辑形式语言所发展出来的方法，在逻辑学、计算机科学和语言学等领域都有广泛应用。所以，从长远来看，仔细钻研这些细节是值得的。

\end{document}